Zum Beweise bedenken wir, da\ss\ eine durch Faltung einer Form $f$ mit einer zweiten $g$ entstandene Form $h$ sich wegen (45.) auffassen l\"a\ss t als eine mehrere kogrediente Variabelnreihen $p_\rho,p'_\rho$ enthaltende Form $f$. Solche Formen lassen sich aber nach der allgemeinen Theorie der Reihenentwicklung darstellen als Polaren der Elementarkovarianten von $f$, dadurch da\ss\ man die Reihenentwicklung sukzessive auf je zwei kogrediente Variabelnreihen $p_\rho,p'_\rho$ anwendet. Dabei sind nach (38.) die auftretenden Elementarkovarianten die durch kogrediente Faltung erzeugten Formen. Da aber die Reihenfolge in der Anwendung der Reihenentwicklung noch willk\"urlich ist, k\"onnen wir insbesondere eine solche Reihenfolge nehmen, die der Erzeugung der allgemeinen Faltungen aus Grundfaltungen entspricht. Das entstehende System von Elementarkovarianten wird dann identisch mit der Formenreihe $f$; und zwar ordnen sich die durch kogrediente Faltung von h\"oherem Defekte entstandenen Formen in die durch Grundfaltungen entstandenen ein. Zu Polaren der Formenreihe $f$ gelangt man aber auch bei beliebiger Reihenfolge der Reihenentwicklung, der ein zweites System von Elementarkovarianten entsprechen m\"oge. Da n\"amlich die Formenreihe die Gesamtheit der invarianten Bildungen ersch\"opft, l\"a\ss t sich jede Form des zweiten Systems linear ausdr\"ucken durch Polaren der Formenreihe, und zwar Polaren derjenigen Formen, die gleiche oder h\"ohere Gewichtsabnahme zeigen. Dies letztere folgt daraus, da\ss\ sich (vgl. \S\ 7) die Polaren auffassen lassen als Simultaninvarianten einer Formenreihe $H$ (52.), mit der Form $h$ entsprechenden Gradzahlen, und der Formenreihe $f$. Jeder Faltung von $H$ in sich entspricht eine Gewichtsabnahme, die sich nach Vornahme von Substitution (11.) auf die Symbolreihen und damit auf die Formen von $f$ \"ubertr\"agt\srcfn{*)}{Diese Betrachtung liegt auch der zweiten Fassung des \"Aquivalenzbegriffes (\S\ 8 Schlu\ss) zugrunde. Weitere Ausf\"uhrungen \"uber die verschiedenen Systeme von Elementarkovarianten finden sich im tern\"aren Gebiet bei Study, a. a. O., II. \S\ 7. Satz 7) und 8). Verm\"oge unseres Satz VII gelten dieselben auch f\"ur $n$ Variable.}.
